\documentclass[12pt,letterpaper]{extarticle}
\usepackage{fontspec}
\usepackage[left=1in,right=1in,top=1in,bottom=1in]{geometry}
\usepackage{tikz}
\usetikzlibrary{calc}
\usepackage{fancyhdr}
\usepackage{enumitem}
\usepackage{microtype}
\usepackage{setspace}
\usepackage{array}
\usepackage{xcolor}
\usepackage{titlesec}
\usepackage{etoolbox}
\usepackage{needspace}
\usepackage{amssymb}
\usepackage[
pdftitle={State v. Penn -- Case No. 24CR49043 -- Brady Giglio Violation Notice},
pdfauthor={Ashle' Penn, Pro Se},
colorlinks=true,linkcolor=blue!70!black,urlcolor=blue!70!black,
bookmarks=true,unicode=true
]{hyperref}
\usepackage{bookmark}
\setmainfont{Times New Roman}
\titleformat{\section}{\normalfont\fontsize{10}{12}\bfseries\centering}
  {\thesection}{1.5em}{}[\vspace{-4pt}]
\titleformat{\subsection}{\normalfont\fontsize{10}{12}\bfseries\centering
  \hyphenpenalty=50\tolerance=9999}{\thesubsection}{1.5em}{}[\vspace{-2.5pt}]
\titlespacing*{\section}{2pt}{6pt}{4pt}
\titlespacing*{\subsection}{2pt}{4pt}{2pt}
\renewcommand{\thesection}{\Roman{section}}
\renewcommand{\thesubsection}{\Alph{subsection}}
\newlength{\PleadingStep}
\setlength{\PleadingStep}{\dimexpr\textheight/28\relax}
\makeatletter
\AtBeginDocument{\setlength{\baselineskip}{\PleadingStep}}
\makeatother
\AddToHook{shipout/background}{%
\begin{tikzpicture}[remember picture,overlay]
\foreach \i in {1,...,28}{
\pgfmathsetmacro{\kk}{\i-1}
\node[anchor=east,font=\fontsize{10pt}{10pt}\selectfont] at
([xshift=0.95in,yshift={-(1in+0.5pt+(\kk+0.5)*\PleadingStep)}]
current page.north west) {\i};}
\draw[line width=0.65pt]
([xshift=0.95in,yshift={-(1in+0.5pt)}]current page.north west)--
([xshift=0.95in,yshift={-(1in+0.5pt+\textheight)}]current page.north west);
\draw[line width=0.65pt]
([xshift=0.975in,yshift={-(1in+0.5pt)}]current page.north west)--
([xshift=0.975in,yshift={-(1in+0.5pt+\textheight)}]current page.north west);
\end{tikzpicture}}
\setlength{\parindent}{0.5in}
\setlength{\parskip}{0pt}
\hyphenpenalty=500\tolerance=2000\emergencystretch=5em\hfuzz=2pt
\raggedbottom
\pagestyle{fancy}\fancyhf{}
\renewcommand{\headrulewidth}{0pt}
\renewcommand{\footrulewidth}{0.4pt}
\setlength{\headheight}{13.6pt}
\fancyfoot[L]{\fontsize{8}{8}\selectfont\textbf{\textsc{%
State v.\ Penn --- Case No.\ 24CR49043 --- Brady/Giglio Violation Notice and Demand%
}} --- Pg.~\thepage}
\newcommand{\CaptionBlock}[3]{%
\noindent\begin{minipage}[t]{3.20in}\vspace{12pt}%
\begin{tikzpicture}
\node[inner sep=0pt,text width=3.20in,anchor=north west] (T) {#3};
\draw[line width=0.5pt](T.south west)--(T.south east);
\draw[line width=0.5pt](T.south east)--(T.north east);
\end{tikzpicture}\end{minipage}%
\hfill\begin{minipage}[t]{3.05in}\vspace{2pt}%
\raggedright Case No.: #1\\[10pt]
\raggedleft\small\textbf{#2}\end{minipage}}
\begin{document}

\begin{center}
{\fontsize{12}{14}\bfseries\selectfont
IN THE CIRCUIT COURT OF THE STATE OF OREGON}\\[4pt]
{\fontsize{12}{14}\bfseries\selectfont
FOR THE COUNTY OF WASHINGTON}\\[18pt]
\end{center}

\CaptionBlock{24CR49043}{%
BRADY AND GIGLIO VIOLATION NOTICE\\[4pt]
AND DEMAND FOR DISCLOSURE\\[8pt]
\normalfont\itshape Special Appearance Only;\\
Filed on the Papers%
}{%
{\bfseries STATE OF OREGON,}\\[6pt]
\hspace{0.3in} Plaintiff,\\[8pt]
v.\\[8pt]
{\bfseries ASHLE' ANTOINETTE PENN,}\\[6pt]
\hspace{0.3in} Defendant.%
}

\vspace{18pt}

\section{PRELIMINARY STATEMENT}

Defendant Ashle' Antoinette Penn appears specially and pro se to provide formal
notice of Brady and Giglio violations in the above-captioned matter and to demand
immediate disclosure of specific, identified materials. This Notice is filed
pursuant to \textit{Brady v.\ Maryland}, 373 U.S.\ 83 (1963); \textit{Giglio v.\
United States}, 405 U.S.\ 150 (1972); \textit{Napue v.\ Illinois}, 360 U.S.\ 264
(1959); ORS 135.815; and Or.\ Const.\ Art.\ I, \S\S~9, 11.

The violations identified herein are not speculative. Each is established by
reference to documents already in the State's possession --- in most instances,
documents authored by the State's own witnesses or agents. Each item of undisclosed
material identified in Section III below is specifically described, its Brady or
Giglio significance is stated, and the basis for believing it exists and has not
been produced is set forth.

Defendant demands production of all items identified in Section III within
\textbf{fourteen (14) days} of service of this Notice. Failure to produce shall be
raised before this Court as a basis for dismissal with prejudice and shall be
preserved as an independent ground for relief on appeal and in any subsequent
civil rights proceeding.

\section{LEGAL FRAMEWORK}

\subsection{The Brady Obligation}

The prosecution has an affirmative constitutional duty to disclose all evidence
favorable to the accused that is material to guilt or punishment. \textit{Brady v.\
Maryland}, 373 U.S.\ 83, 87 (1963). Evidence is material if there is a reasonable
probability that its disclosure would have produced a different result. \textit{United
States v.\ Bagley}, 473 U.S.\ 667, 682 (1985). The duty attaches to evidence
possessed by any member of the prosecution team, including law enforcement.
\textit{Kyles v.\ Whitley}, 514 U.S.\ 419, 437 (1995).

Under Oregon law, ORS 135.815 requires the State to disclose, upon request, all
relevant written or recorded statements of witnesses, all reports of scientific
tests, and all evidence known to the district attorney that tends to negate the
guilt of the accused or mitigate the offense. This obligation is independent of
and in addition to the constitutional Brady requirement.

\subsection{The Giglio/Napue Obligation}

The prosecution must disclose evidence bearing on the credibility of its own
witnesses. \textit{Giglio v.\ United States}, 405 U.S.\ 150, 154 (1972). The
prosecution may not present testimony it knows to be false and may not allow false
testimony to stand uncorrected. \textit{Napue v.\ Illinois}, 360 U.S.\ 264, 269
(1959). A conviction obtained through false testimony violates due process
regardless of the prosecution's good faith. \textit{Id.}

\subsection{The Continuing Duty}

The Brady obligation is continuing. \textit{Strickler v.\ Greene}, 527 U.S.\ 263,
280 (1999). The State's obligation to disclose does not cease with indictment. It
extends through trial. The State is in violation of that duty at every moment an
identified material item remains undisclosed.

\section{IDENTIFIED BRADY AND GIGLIO VIOLATIONS}

\subsection{Category One: Rambin's December 27, 2024 Supplemental Report and
the Misidentification Admission}

\textbf{The Material:} Deputy David Rambin (DPSST \#53879) filed a supplemental
report on December 27, 2024 in case file GO 50 2024-13459 stating that it was
possible Defendant was not the suspect in the September 17, 2024 incident. The
same report documents that Rambin returned to Target on December 27, 2024 and
directed Target loss prevention to compile all available information on Defendant
separately, as though building a new case from scratch.

\textbf{Brady Significance:} This report is the arresting officer's own admission
that the identification predicate for the September 18, 2024 arrest may have been
wrong. It is exculpatory on its face. It directly undermines the probable cause
basis for both counts of the January 14, 2025 indictment. A reasonable probability
exists that its disclosure before the grand jury proceeding would have produced a
different result.

\textbf{Status:} This report was not disclosed to Defendant before the January 14,
2025 grand jury proceeding. The grand jury convened seventeen days after the report
was filed. The State presented the case to the grand jury without disclosing that
the arresting officer had placed the identification in doubt. The omission was not
inadvertent --- the State excluded from the grand jury presentation the very facts
(the September 17 incident and its associated valuation) that Rambin's report
placed in question, demonstrating that the State knew of the report and structured
the grand jury presentation around it.

\textbf{Demand:} Produce the complete December 27, 2024 supplemental report,
all attachments, all communications transmitting the report to the DA's office,
and all records reflecting the date on which the DA's office received the report.

\subsection{Category Two: Body Worn Camera Footage --- September 18, 2024
Evidence Pickup at Target (Approximately 7:00 PM)}

\textbf{The Material:} Deputy Rambin's body worn camera footage from his
approximately 7:00 PM visit to Target on September 18, 2024, during which he
retrieved the thumb drive containing September 17 surveillance footage. This is a
separate and distinct recording from the arrest BWC footage.

\textbf{Brady Significance:} This footage documents the entirety of Rambin's
interaction with Target personnel during the evidence pickup. Specifically, it
establishes: (a) that Kinley Baker, the LP officer whose name appears in the
grand jury testimony, was not present during the evidence pickup; (b) that the
thumb drive was retrieved from an unsecured front desk location, not through any
chain-of-custody protocol; (c) that Rambin's only communication with Sarah Fillis
consisted of a brief exchange from a distance while Fillis was in transit from
lunch, lasting a matter of seconds; and (d) that no conversation occurred between
Rambin and Fillis containing the detailed description of the September 17 incident
that Rambin's probable cause affidavit attributes to Fillis as a verbal statement.

The affidavit states: ``When I arrived, Sarah stated that a black female adult had
entered the store and after selecting multiple items, left the store through a fire
exit without paying for the items on September 17, 2024.'' The BWC footage of
that visit, if produced, will establish whether any such conversation occurred.
Defendant submits it did not.

\textbf{Status:} Not produced in discovery. Its existence is established by
WCSO's own BWC template documentation filed in the case. The footage is in the
possession of WCSO.

\textbf{Demand:} Produce the complete BWC footage from Deputy Rambin's
September 18, 2024 visit to Target at approximately 7:00 PM, including all
pre-activation buffer footage, and all CAD entries reflecting Rambin's location
and activity between approximately 7:00 PM and 9:42 PM (2142hrs) on September 18,
2024.

\subsection{Category Three: CAD Records Reflecting Rambin's Location and
Activity Between the Evidence Pickup and the Arrest}

\textbf{The Material:} Computer-Aided Dispatch records for Deputy Rambin
(DPSST \#53879) from approximately 7:00 PM through 9:42 PM on September 18, 2024,
reflecting his location, status, and documented activity during the approximately
two hours and forty minutes between the evidence pickup and the arrest.

\textbf{Brady Significance:} The probable cause affidavit implicitly represents
that Rambin reviewed the September 17 surveillance footage before arresting
Defendant. His own primary follow-up narrative, submitted September 19, 2024,
states that he ``was later able to review the security report'' and learned details
of the September 17 incident from that review. The question of \textit{when} that
review occurred --- before or after the arrest --- is dispositive. If Rambin
watched the footage after the arrest, the detailed description of the September 17
incident in the affidavit was not derived from pre-arrest knowledge and was
attributed to Sarah Fillis falsely.

CAD records will establish where Rambin was and what he was doing during the
window between the evidence pickup and the arrest. If those records show Rambin
drove directly from Target to Bethany Station and began processing Defendant's
booking, they establish that the footage review was post-arrest. That determination
transforms an inconsistency into a material misstatement in a sworn affidavit.

\textbf{Status:} Not produced in discovery.

\textbf{Demand:} Produce all CAD entries, status logs, dispatch records, and
mobile data terminal records for Deputy Rambin from 6:30 PM through midnight on
September 18, 2024.

\subsection{Category Four: Communications Between the DA's Office and WCSO
Between September 18, 2024 and January 14, 2025}

\textbf{The Material:} All written communications --- including emails, text
messages, case management system entries, and interoffice memoranda --- between
the Washington County District Attorney's office and any member of the Washington
County Sheriff's Office regarding case 24CR49043 and/or the investigation of
Defendant, from the date of arrest through the date of the grand jury proceeding.

\textbf{Brady/Giglio Significance:} These communications will establish: (a) the
date on which the DA's office received Rambin's December 27, 2024 supplemental
admission; (b) whether any member of the prosecution team discussed the
misidentification problem with any member of WCSO prior to the grand jury; (c)
whether the decision to exclude the September 17 incident from the grand jury
presentation was deliberate and coordinated; and (d) whether any member of the
prosecution team knew that Rambin's and Baker's grand jury testimony conflicted
with the sworn September 18 affidavit and chose not to correct it.

The significance of this category is categorical: if the communications show the
DA's office received the December 27 supplemental report before January 14 and
proceeded to the grand jury without disclosure, the \textit{Napue} violation is
established from the State's own records.

\textbf{Status:} Not produced. Not disclosed.

\textbf{Demand:} Produce all written communications between the Washington County
District Attorney's office and WCSO, including Deputy Rambin, regarding case
24CR49043, Defendant Ashle' Penn, and/or the September 18, 2024 arrest, from
September 18, 2024 through January 14, 2025.

\subsection{Category Five: The Grand Jury Transcript}

\textbf{The Material:} The complete transcript of the January 14, 2025 grand jury
proceeding in case 24CR49043, including the testimony of Deputy David Rambin and
Kinley Baker, all exhibits admitted, and all statements by the prosecutor to the
grand jury.

\textbf{Brady/Giglio/Napue Significance:} The transcript is necessary to establish
the \textit{Napue} violation on the record. The demurrer filed April 28, 2025
identifies a specific, documentable conflict: Rambin's September 18, 2024
affidavit attributes the initial tip to Sarah Fillis contacting Deputy Orozco.
Baker's grand jury testimony attributes the initial contact to Baker contacting
Rambin directly. These accounts are mutually exclusive on a material fact. The
prosecutor had the affidavit in the case file. The transcript will establish
whether the prosecutor corrected Baker's testimony, presented the conflict to the
grand jury, or allowed it to stand without correction.

Additionally, the transcript will establish whether the grand jury was told: (a)
that Rambin's supplemental report admitted possible misidentification; (b) that the
September 17 facts were being excluded because of that admission; and (c) that the
investigation of the September 18 conduct was initiated by a warrantless arrest
whose probable cause basis had been placed in doubt by the arresting officer
himself.

\textbf{Standard for Disclosure:} Grand jury secrecy under ORS 132.210 yields to
defendant's particularized need for the transcript when the defendant makes a
showing that the testimony is likely to conflict with a sworn instrument in the
State's file and that the conflict is material to a pending constitutional
challenge. \textit{State v.\ Farber}, 314 Or.\ 397 (1992). That showing is made
here. The conflict between the affidavit and Baker's testimony is identified
specifically. The pending demurrer and motion to dismiss are the proceedings for
which the transcript is needed. The standard is satisfied.

\textbf{Demand:} Produce the complete January 14, 2025 grand jury transcript,
or, in the alternative, submit it to the Court for in camera review with a
proposed order directing production of the portions reflecting the testimony
of Rambin and Baker and all prosecutor statements regarding the probable cause
basis for the charges.

\subsection{Category Six: The Original Email from Sheila Clark to Judge Thompson
and All Related Communications}

\textbf{The Material:} The complete text of the email sent by Sheila Clark,
P\&P Services Supervisor at Washington County Community Justice, to the Honorable
Brandon Thompson, on or about June 10, 2025, regarding Defendant, including all
prior correspondence in the same thread, all attachments, and all subsequent
communications between Judge Thompson and any member of Washington County Community
Justice or the Washington County Circuit Court Pretrial Release Office regarding
Defendant.

\textbf{Brady Significance:} The Stiefer email chain already in the record
establishes that Clark sent an email to Thompson and that Thompson forwarded it to
the Pretrial Release Office and directed show cause proceedings. The content of
Clark's original email is unknown. That content is material because: (a) it is the
initiating instrument for a proceeding that resulted in an arrest warrant; (b) it
reflects what Clark actually knew or believed at the time she contacted the judge,
before any investigation had been conducted; (c) it may reveal the source of
Clark's information about Defendant's release conditions; and (d) it goes directly
to whether the show cause proceeding was initiated on a factual basis or on
institutional animus against Defendant arising from her advocacy on behalf of the
co-defendant in his supervision proceedings.

\textbf{Status:} Never produced. Never disclosed. The email exists in Thompson's
account, Clark's account, and potentially in the pretrial release office's records.
Its non-production while the proceeding it initiated was being conducted against
Defendant is itself a Brady violation.

\textbf{Demand:} Produce the complete text of all communications between Sheila
Clark and Judge Brandon Thompson regarding Defendant Ashle' Penn, including the
email Clark sent to Thompson on or about June 10, 2025, all prior communications
in the same thread, all forwarded copies, and all communications between Thompson's
chambers and the Washington County Circuit Court Pretrial Release Office regarding
Defendant, from January 1, 2025 through the date of this Notice.

\subsection{Category Seven: All Records Reflecting How Clark Obtained Information
About Defendant's Release Conditions}

\textbf{The Material:} All records --- including case management system entries,
inter-agency communications, shared database access logs, and written reports ---
reflecting how Sheila Clark or any member of Washington County Community Justice
obtained information about Defendant's release conditions in case 24CR49043 prior
to June 10, 2025.

\textbf{Brady Significance:} Clark's agency supervised the co-defendant in a
separate matter. Defendant's release agreement is a document in her criminal case
file. The mechanism by which Clark obtained information about Defendant's release
conditions is material to whether that information was obtained through proper
channels or through unauthorized access to Defendant's case file, and to whether
the show cause proceeding was initiated on the basis of legitimate compliance
information or on information gathered and used outside the scope of Clark's
supervisory authority.

\textbf{Demand:} Produce all records reflecting how Washington County Community
Justice personnel obtained knowledge of Defendant's release conditions, including
any inter-agency information sharing agreements, shared database access logs for
Defendant's case file, and written communications between Washington County
Community Justice and Washington County Circuit Court personnel regarding Defendant
prior to June 10, 2025.

\subsection{Category Eight: Target Surveillance Footage --- September 17, 2024}

\textbf{The Material:} All surveillance footage from the Target store at 12675 NW
Cornell Road, Portland, Oregon, from September 17, 2024, including footage of the
person identified by loss prevention as having exited through a fire door without
paying for merchandise.

\textbf{Brady Significance:} The September 17 surveillance footage is the
foundational document of this prosecution. The September 18, 2024 probable cause
affidavit attributes to Defendant conduct that occurred on September 17. Rambin's
December 27, 2024 supplemental report acknowledges a possible misidentification ---
and the most direct evidence bearing on that misidentification is the footage of
the actual September 17 suspect. That footage will either corroborate or refute the
identification. If the September 17 footage shows a person whose physical
characteristics differ materially from Defendant's, the misidentification is
established visually. That is exculpatory evidence of the highest order.

\textbf{Status:} Not produced. The thumb drive Rambin retrieved on September 18,
2024 reportedly contained this footage. The chain of custody of that thumb drive ---
retrieved from an unsecured desk, confirmed from a distance by Fillis, without
Baker present --- is itself a matter of record. The footage and the chain of custody
records must both be produced.

\textbf{Demand:} Produce all surveillance footage from Target at 12675 NW Cornell
Road from September 17, 2024, all chain of custody documentation for the thumb
drive retrieved by Rambin on September 18, 2024, and all records reflecting the
disposition of that footage.

\subsection{Category Nine: The Barry Washington Trial Transcript, Pages 184--186}

\textbf{The Material:} Pages 184 through 186 of the trial transcript in the matter
of \textit{State v.\ Barry Washington}, Washington County Circuit Court case
24CR49041, reflecting the testimony of Sarah Fillis.

\textbf{Brady/Giglio Significance:} Sarah Fillis is identified in Rambin's
September 18, 2024 probable cause affidavit as the source of the identification
that led to Defendant's arrest. At the Washington trial, Fillis testified in a
manner that contradicts the affidavit's account of her communications with law
enforcement. Specifically, Fillis's testimony at pages 184--186 of the Washington
trial transcript addresses who she identified to deputies and in what terms.
That testimony is exculpatory because it directly contradicts the sworn instrument
that forms the basis for this prosecution. It is Giglio material because it
impeaches the foundational reliability of Rambin's affidavit. It triggers a
continuing Brady disclosure obligation because it became available to the State at
the moment Fillis testified under oath in a related case in the same courthouse.

\textbf{Status:} Not produced. The State was present at the Washington trial. The
continuing Brady obligation attached at the moment Fillis testified. The obligation
has not been discharged.

\textbf{Demand:} Produce pages 184 through 186 of the \textit{State v.\ Barry
Washington} trial transcript, case 24CR49041, reflecting the testimony of Sarah
Fillis, and all notes, memoranda, or summaries of that testimony in the possession
of any member of the prosecution team.

\section{NOTICE OF CONSEQUENCES OF NON-COMPLIANCE}

The items demanded in Section III are not marginal or peripheral. Each is either:
(a) a document authored by the State's own witness that contradicts a sworn
instrument in the State's file; (b) a recording that will establish the factual
basis --- or the absence of a factual basis --- for representations made in that
sworn instrument; or (c) communications between government actors that will
establish what the prosecution knew and when.

The State's failure to produce these materials within fourteen days of this Notice
will be presented to this Court as: (1) an independent basis for dismissal with
prejudice under \textit{Brady}, ORS 135.815, and Or.\ Const.\ Art.\ I, \S~11;
(2) an aggravating factor in support of the pending motion to dismiss for
prosecutorial misconduct; and (3) grounds for an adverse inference instruction
should the case proceed to trial.

Defendant further places the State on notice that the non-disclosure of these
materials has already caused prejudice. The January 14, 2025 grand jury proceeded
without knowing what the December 27, 2024 supplemental report said. The demurrer
has been pending for fourteen months in part because the materials needed to fully
develop the constitutional challenge remain in the State's exclusive possession.
Every day of continued non-disclosure extends the prejudice.

\section{REQUEST FOR JUDICIAL RELIEF}

Defendant respectfully requests that this Court:

\begin{enumerate}[leftmargin=0.5in, label=\arabic*., itemsep=4pt]

\item Take judicial notice of this Notice and Demand and enter it in the record
of case 24CR49043.

\item Order the State to produce all materials identified in Section III within
fourteen (14) days of this Court's order.

\item Order the State to submit to the Court, within fourteen (14) days, a written
declaration identifying any item listed in Section III that the State contends does
not exist, has been destroyed, or is not in the prosecution team's possession,
with a statement of when and how each such item was lost or destroyed.

\item Set a status conference within thirty (30) days to address compliance with
this Order and to schedule argument on the pending demurrer and the motions
identified in Defendant's concurrently filed Notice of Appearance.

\end{enumerate}

\vspace{24pt}
\noindent\hspace{0.5in}Respectfully submitted,

\vspace{48pt}
\noindent\hspace{0.5in}\rule{2.5in}{0.4pt}\\[4pt]
\noindent\hspace{0.5in}Ashle' Antoinette Penn\\
\noindent\hspace{0.5in}Defendant, Pro Se\\[6pt]
\noindent\hspace{0.5in}Dated: \today\\[6pt]
\noindent\hspace{0.5in}\textit{Address:} \underline{\hspace{2.2in}}\\[4pt]
\noindent\hspace{0.5in}\textit{Phone:} \underline{\hspace{2.2in}}\\[4pt]
\noindent\hspace{0.5in}\textit{Email:} \underline{\hspace{2.2in}}

\newpage
\begin{center}{\fontsize{11}{13}\bfseries\selectfont CERTIFICATE OF SERVICE}
\end{center}
\vspace{12pt}
I, Ashle' Antoinette Penn, certify that on \today\ I served a true and correct
copy of the foregoing \textit{Brady and Giglio Violation Notice and Demand for
Disclosure} upon the Washington County District Attorney's Office,
150 N.\ First Avenue, Suite 300, Hillsboro, Oregon 97124, by the method indicated:
\vspace{12pt}
\noindent\hspace{0.5in}$\square$\ \ Electronic Filing (eFiling)
\hspace{0.8in}
$\square$\ \ U.S.\ Mail, First Class
\vspace{36pt}
\noindent\hspace{0.5in}\rule{2.5in}{0.4pt}\\[4pt]
\noindent\hspace{0.5in}Ashle' Antoinette Penn

\end{document}
